% Copyright 2026 The terCAD team. All rights reserved.
% Use of this content is governed by a CC BY-NC-ND 4.0 license that can be found in the LICENSE file.

\subsubsection{Apple Store - iOS} \label{d-ios}

The iOS application is an Xcode project (\q{iosApp/module.xcodeproj}), and its "Build Kotlin" phase compiles the shared
code into the \q{KotlinModules}-framework by Gradle (\q{./gradlew :iosApp:embedAndSignAppleFrameworkForXcode}). So that, there is no need in a device running macOS:
the whole build, signing and upload is made by GitHub Actions on a \q{macos-26} runner (\q{iOS}-target of the \q{build}-matrix in
\q{.github/workflows/build.yml}). All that is left to be done manually is to get the signing assets from Apple, register
the application and put the assets into GitHub secrets.

\paragraph{Prerequisites.} An Apple Developer Program membership (organization) and access to
\href{https://appstoreconnect.apple.com/}{https://appstoreconnect.apple.com/} with "Admin" or "Account Holder" role. The
\q{Team ID} can be taken from \href{https://developer.apple.com/account}{https://developer.apple.com/account} ("Membership
details"). Next values are already defined in the repository and have to be consistent with those that are registered
by Apple:

\begin{lstlisting}[language=terminal]
Bundle ID:         com.tercad.zwyka   # iosApp/module.xcodeproj/project.pbxproj
Profile name:      Zwyka iOS          # iosApp/ExportOptions.plist, build.yml (PROFILE_NAME)
Version:           <tag without "v">  # MARKETING_VERSION, e.g. 1.2.0 (max three integers)
Build number:      <commits count>    # CURRENT_PROJECT_VERSION, has to increase on each upload
\end{lstlisting}

\paragraph{Step 1: Distribution certificate.} The "Apple Distribution" certificate is issued per team, so that the one
that is already used by other applications of terCAD (secrets \q{APPLE\_P12} and \q{APPLE\_P12\_KEY}) can be reused, and
this step skipped. Otherwise, it's needed to create a private key and a certificate signing request (by \q{openssl}, or
via "Keychain Access" on macOS):

\begin{lstlisting}[language=terminal]
# Create a private key
$ openssl genrsa -out apple.key 2048
# Create the CSR file (C - Country Name [2 letter code])
$ openssl req -new -key apple.key \
    -out CertificateSigningRequest.certSigningRequest \
    -subj "/emailAddress=your@mail.to, CN=Your Name, C=XX"
\end{lstlisting}

\noindent and upload it to
\href{https://developer.apple.com/account/resources/certificates/list}{https://developer.apple.com/account/resources/certificates/list}
by choosing "\emph{[Software] Apple Distribution: Sign your iOS, iPadOS, macOS, tvOS, watchOS, and visionOS apps for
release testing using Ad Hoc distribution or for submission to the App Store}". Downloaded \q{.cer}-file has to be
converted into \q{.p12}, and then into a string (base64) representation:

\begin{lstlisting}[language=terminal]
# Convert the .cer-file into .pem-format
$ openssl x509 -in distribution.cer -inform DER -out apple.pem -outform PEM
# Use .pem-file and private .key to generate .p12
$ openssl pkcs12 -export -legacy -out apple.p12 -inkey apple.key -in apple.pem
Enter Export Password: #####
# Convert file to base64 string
$ base64 apple.p12 > apple.txt
\end{lstlisting}

\noindent Usage of \q{legacy}-attribute is a key factor to avoid a security failure "SecKeychainItemImport: MAC
verification failed" since OpenSSL 3.x has changed its default algorithm in \q{pkcs12} which is not compatible with
embedded Security frameworks in macOS/iOS. Private key and \q{.p12}-file have to be stored in a safe place, outside the
repository.

\paragraph{Step 2: Identifier.} Register an App ID on
\href{https://developer.apple.com/account/resources/identifiers/list}{https://developer.apple.com/account/resources/identifiers/list}:
"\emph{App IDs}" $\rightarrow$ "\emph{App}", description "Zwyka", explicit Bundle ID \q{com.tercad.zwyka}. No additional
capabilities are required: reminders are local notifications, and Siri Shortcuts are provided by App Intents (both
of them don't need entitlements).

\paragraph{Step 3: Provisioning profile.} On
\href{https://developer.apple.com/account/resources/profiles/list}{https://developer.apple.com/account/resources/profiles/list}
create "\emph{Distribution} $\rightarrow$ \emph{App Store Connect}" profile for the App ID from the previous step, select
the distribution certificate, and name it exactly "\emph{Zwyka iOS}". The downloaded file is converted into base64:

\begin{lstlisting}[language=terminal]
$ base64 Zwyka_iOS.mobileprovision > apple-provision.txt
\end{lstlisting}

\noindent The profile expires together with the certificate (in a year): then it has to be regenerated, and the secret
replaced.

\paragraph{Step 4: Application record.} On \href{https://appstoreconnect.apple.com/apps}{https://appstoreconnect.apple.com/apps}
"\emph{New App}": platform iOS, name "Zwyka", primary language English (U.S.), Bundle ID \q{com.tercad.zwyka}, SKU
(any unique value, e.g. \q{zwyka}), full access. Name has to be unique across the App Store, otherwise a subtitle-like
variant is required (e.g. "Zwyka: Habit Tracker").

\paragraph{Step 5: API key for uploading.} On
\href{https://appstoreconnect.apple.com/access/integrations/api}{https://appstoreconnect.apple.com/access/integrations/api}
generate a "Team Key" with "App Manager" access (one key can be shared across applications of the team). Copy
\q{Issuer ID} and \q{KEY ID}, download \q{AuthKey\_<KEY ID>.p8} (it can be downloaded only once) and convert it:

\begin{lstlisting}[language=terminal]
$ base64 -w 0 AuthKey_XXXXXXXXXX.p8 > apple-api.txt
\end{lstlisting}

\paragraph{Step 6: GitHub secrets.} Add next values to "Settings $\rightarrow$ Secrets and variables $\rightarrow$
Actions" of the repository (or to the organization secrets, to be shared with other applications):

\begin{lstlisting}[language=terminal]
APPLE_P12          # content of apple.txt (Step 1)
APPLE_P12_KEY      # export password of apple.p12
APPLE_IOS_PROFILE  # content of apple-provision.txt (Step 3)
APPLE_TEAM_ID      # Team ID, e.g. 8PVKPQ758L
APPLE_API_P8       # content of apple-api.txt (Step 5)
APPLE_API_KEY      # KEY ID (Step 5)
APPLE_API_ISSUE    # Issuer ID (Step 5)
\end{lstlisting}

\noindent After that, each pushed tag \q{v*} archives the application by \q{xcodebuild archive} with manual signing
(the signing parameters are passed as build settings, so that the local development with automatic signing is not
affected), exports \q{zwyka\_iOS.ipa} by \q{iosApp/ExportOptions.plist}, attaches it to the release page, and uploads to
App Store Connect by \q{xcrun altool --upload-app}. The pipeline ends with \q{UPLOAD SUCCEEDED}, and in a few minutes
the build appears in the "TestFlight"-tab.

\paragraph{Step 7: Compliance.} \q{ITSAppUsesNonExemptEncryption} is set to \q{false} in \q{iosApp/src/Info.plist}
(no network communication, and no encryption beyond the operating system), so the "Export Compliance Information"
questions are not asked for each build. It has to be revised if the application starts using encryption (for example, to
sync data): then the flag is set to \q{true}, and \q{ITSEncryptionExportComplianceCode} is added from
\href{https://appstoreconnect.apple.com/apps}{appstoreconnect.apple.com/apps/\q{id}/features/encryption}.

The privacy manifest \q{iosApp/src/PrivacyInfo.xcprivacy} declares the usage of "required reason" APIs (UserDefaults
for the chosen language; file timestamp, system boot time and disk space by Kotlin/Native runtime, SQLite and Compose
Multiplatform), otherwise the upload is rejected with "ITMS-91053: Missing API declaration". The manifest has to be
extended if new APIs of that kind are used.

\paragraph{Step 8: Store listing.} On the application page (\q{App Store}-tab) it's needed to fill in:

\begin{itemize}
  \item "\emph{App Information}": subtitle, category (Health \& Fitness, or Productivity), content rights, age rating
  questionnaire (no objectionable content $\rightarrow$ 4+);
  \item "\emph{App Privacy}": privacy policy URL, and "\emph{Data Not Collected}" (all the data is kept on the device);
  \item "\emph{Pricing and Availability}": price (free) and countries;
  \item version page: description, keywords, support URL, promotional text, screenshots (6.9" iPhone -- 1320x2868 or
  1290x2796, and 13" iPad -- 2064x2752 or 2048x2732, since iPad is supported by \q{TARGETED\_DEVICE\_FAMILY}), and
  localizations for the supported languages;
  \item "\emph{Business}" (Agreements, Tax, and Banking): Free Apps Agreement is accepted, and the trader status for the
  European Union (Digital Services Act) is declared.
\end{itemize}

\paragraph{Step 9: Review.} Select the uploaded build in the "\emph{Build}"-section of the version page, and "\emph{Add
for Review}" $\rightarrow$ "\emph{Submit for Review}". Next versions repeat only this step: create a new version (equal
to the tag), fill "What's New in This Version", select the build and submit.

\paragraph{TestFlight (optional).} Pre-release builds can be available for external testing: create an "\emph{External
Testing}"-group in the "TestFlight"-tab, add a build to it and submit it for the beta review. After the approval,
"\emph{Enable Public Link}" gives the link for testers.
